\corrigesChapitre{2}{Module, arguments et plan complexe}
\begin{corrige}{1}{Affixes}
\begin{questions}
\item $-2+3\mathrm i$.
\item $(4,-1)$.
\item $b-a=2+3\mathrm i$.
\end{questions}
\end{corrige}
\begin{corrige}{2}{Modules}
\begin{questions}
\item $\sqrt{9+16}=5$.
\item $2$ et $5$.
\item $|1+\mathrm i|^5=(\sqrt2)^5=4\sqrt2$.
\end{questions}
\end{corrige}
\begin{corrige}{3}{Arguments remarquables}
\begin{questions}
\item $2k\pi$ pour $1$ ; $\pi/2+2k\pi$ pour $\mathrm i$, avec $k\in\Z$.
\item $\pi+2k\pi$, $k\in\Z$.
\item Le vecteur nul ne définit pas de direction ; la forme trigonométrique ne sélectionne aucun angle.
\end{questions}
\end{corrige}
\begin{corrige}{4}{Forme trigonométrique}
\begin{questions}
\item $\sqrt2(\cos(\pi/4)+\mathrm i\sin(\pi/4))$.
\item $2(\cos(5\pi/6)+\mathrm i\sin(5\pi/6))$.
\item $\sqrt3+\mathrm i$.
\end{questions}
\end{corrige}
\begin{corrige}{5}{Distances et milieu}
\begin{questions}
\item $|4-3\mathrm i|=5$.
\item $(a+b)/2=3+\mathrm i/2$.
\item $2m-a=b=5-\mathrm i$, comme attendu.
\end{questions}
\end{corrige}
\begin{corrige}{6}{Le cercle unité}
\begin{questions}
\item $|z|=5/5=1$.
\item $1/z=\bar z=(3-4\mathrm i)/5$.
\item $|z^{10}|=|z|^{10}=1$.
\end{questions}
\end{corrige}
\begin{corrige}{7}{Décrire des lieux}
\begin{questions}
\item Cercle de centre d’affixe $2-\mathrm i$ et de rayon $3$.
\item Disque fermé de centre $O$ et de rayon $2$.
\item Médiatrice des points d’affixes $1$ et $-1$, donc axe imaginaire.
\end{questions}
\end{corrige}
\begin{corrige}{8}{Un argument mal choisi}
\begin{questions}
\item $\sqrt2$.
\item $-3\pi/4$, ou $5\pi/4$ modulo $2\pi$.
\item La tangente seule ne distingue pas deux angles séparés de $\pi$. Les signes du cosinus et du sinus doivent être tous deux négatifs.
\end{questions}
\end{corrige}
\begin{corrige}{9}{Démontrer le produit}
\begin{questions}
\item $|zw|^2=zw\overline{zw}=zw\bar z\bar w$.
\item Par commutativité, $(z\bar z)(w\bar w)=|z|^2|w|^2$.
\item Les deux modules sont positifs ou nuls, donc l’unique racine positive est $|z||w|$.
\end{questions}
\end{corrige}
\begin{corrige}{10}{Médiatrice par le calcul}
\begin{questions}
\item $\sqrt{(x+1)^2+y^2}=\sqrt{(x-3)^2+y^2}$.
\item Élever au carré donne $x^2+2x+1=x^2-6x+9$, donc $x=1$.
\item La droite verticale $x=1$ passe par le milieu du segment réel $[-1,3]$ et lui est perpendiculaire.
\end{questions}
\end{corrige}
\begin{corrige}{11}{Triangle et égalité}
\begin{questions}
\item $|3+3\mathrm i|=3\sqrt2$ et $\sqrt2+2\sqrt2=3\sqrt2$.
\item $w=2z$ : les vecteurs ont même direction et même sens.
\item $|z+w|=0$ mais $|z|+|w|=2\sqrt2$ ; l’inégalité est alors stricte.
\end{questions}
\end{corrige}
\begin{corrige}{12}{Module constant ?}
\begin{questions}
\item Sa partie imaginaire est $1$, donc il ne peut être nul.
\item Le numérateur et le dénominateur sont conjugués et ont même module $\sqrt{x^2+1}$.
\item Non : $x-\mathrm i=x+\mathrm i$ imposerait $-\mathrm i=\mathrm i$, impossible. Le lieu est contenu dans le cercle unité privé de ce point.
\end{questions}
\end{corrige}
